\documentclass[10pt]{article}
\usepackage{template/template}
\usepackage{graphicx}
\NewsSetKeys{date={2026年11月4日},vol={022},title={每日有趣新闻},author={ChatGPT},note={核查短刊：截至北京时间11月4日上午；因采编未达到完整发行标准，本稿仅为待扩充源稿。}}
\begin{document}
\begin{documents}
\newsitem{Swift望远镜即将再入大气层：22年高能宇宙观测告一段落}{航天·NASA，11月2日更新}{NASA表示，2004年发射的Neil Gehrels Swift天文台预计于11月4日重返地球大气层。美国太空军在11月2日给出的预测时间为美国东部时间11月4日凌晨3时30分，误差约11小时。该约3200磅的卫星曾以三个望远镜快速追踪伽马射线暴及其他瞬变天体，原计划执行两年，却运行了近22年。太阳活动增强导致高层大气膨胀，加速其轨道衰减。NASA预计大部分结构在再入过程中烧毁，但部分零件可能存留；估计全球有人受伤的概率约为1/2100。该预测具有时间不确定性，不能将其当作已经发生的再入事件。\cite{nasaSwift}}
\newsitem{美国中期选举：民调与选举结果不可混同}{国际政治·11月3日}{11月3日美国中期选举投票进行中。路透社与益普索在选前对1765名美国成年人开展的调查显示，生活成本是选民的重要关切。该调查不是出口民调，更不是官方计票；样本覆盖成年人，不能简单映射为各州或选区席位。国会控制权必须依赖各选区实际投票、计票与认证。选前经济焦虑与通胀压力是可观察的民意线索，不能据此提前宣布任何一党获胜。\cite{reutersPoll}}
\newsitem{数学拾趣：四维Costas阵列反例}{组合数学·arXiv 11月3日}{Shiv Munagala提交预印本，声称构造了阶为6的四维周期Costas阵列，反驳Rubio与Torres关于周期延拓的一项猜想。Costas条件要求不同点对的位移不重复；在周期窗口中维持这一条件是更强的约束。作者称只需检查36个窗口，且已通过计算机及Lean证明助手验证。这里的“已验证”属于作者报告，并不等于本刊独立复现，更不意味着已经通过同行评议。\cite{munagala}}
\newsitem{数学拾趣：Reed--Muller码的覆盖半径}{编码理论·arXiv 11月3日}{Xiang-Dong Hou、Evelyn Sabo与Freeman Slaughter在预印本中声称，一阶九变量Reed--Muller码的覆盖半径精确等于242，改进了原先242至244的已知范围。覆盖半径是所有二进制向量到码中最近码字距离的最大值。作者使用群作用下的轨道分类和计算机辅助检验，并提供复核笔记本。数值及证明仍需独立验证，尚不能作为同行评议后的既定结果。\cite{hou}}
\end{documents}
\bibliographystyle{plain}
\bibliography{refs}
\end{document}
